% ===============================================================
\section{Related Work}\label{sec:related_work}
% ===============================================================

\noindent\textbf{Retrieval-Augmented Generation.}
RAG combines parametric knowledge from neural language models with
non-parametric knowledge retrieved from an external
corpus~\cite{lewis2020rag,guu2020realm}. Dense retrieval is effective for
open-domain factoid QA but degrades on multi-hop evidence
chains~\cite{karpukhin2020dpr,trivedi2023interleaving}. A line of advanced
variants improves retrieval control---Self-RAG through
self-reflection~\cite{asai2023selfrag}, FLARE through active
retrieval~\cite{jiang2023flare}, RAPTOR through hierarchical
summaries~\cite{sarthi2024raptor}, and CRAG through corrective
re-retrieval~\cite{yan2024corrective}. These methods refine \emph{how} dense
retrieval is performed, but they operate within a single evidence mechanism
and do not decide whether a query requires relational traversal or
clarification.

\noindent\textbf{Graph Retrieval and GraphRAG.}
Graph-based retrieval lets evidence follow explicit
relationships~\cite{sun2019pullnet,wang2024kgp}. GraphRAG builds entity
graphs for query-focused summarisation~\cite{edge2024graphrag}; HippoRAG
introduces a neurobiologically inspired graph memory for associative
retrieval~\cite{gutierrez2024hipporag}; LightRAG uses dual-level retrieval
over entity and community knowledge~\cite{he2024lightrag}; and
agent-style methods such as Think-on-Graph~\cite{sun2024thinkongraph} and
GraphReader~\cite{li2024graphreader} let a language model iteratively explore
the graph during reasoning. For legal QA, graph retrieval is especially
relevant because legal documents encode structured relations such as
amendment, replacement, and jurisdiction~\cite{sovrano2020legalkg}. However,
these systems assume graph traversal is always applied; none decides
\emph{when} traversal is warranted versus when a dense lookup would suffice at
lower cost.

\noindent\textbf{Adaptive Routing and Cost-aware Inference.}
Adaptive-RAG selects a retrieval strategy according to question
complexity~\cite{jeong2024adaptive}, while FrugalGPT and RouteLLM route
queries across models of different
cost~\cite{chen2023frugalgpt,ong2024routellm}. These works establish the
cost-aware cascade principle that motivates our design. Their routing target,
however, is the \emph{model} or \emph{retrieval depth}; they do not route
among qualitatively different evidence mechanisms, and they do not treat
clarification as a routable action. Adaptive-RAG is the closest method in
spirit, but it was developed for English open-domain QA and routes among
zero-, single-, and multi-step retrieval over a single dense index, so it is
not directly transferable to a Vietnamese legal corpus that must additionally
choose between dense and graph evidence; we discuss this comparability gap in
Section~\ref{sec:limitations}.

\noindent\textbf{Ambiguity Handling.}
CANARD studies question rewriting for context-dependent
questions~\cite{elgohary2019canard}, and AmbigQA shows that open-domain
questions often admit multiple valid interpretations~\cite{min2020ambigqa}.
These studies treat ambiguity resolution as a stand-alone task applied before
or after retrieval, rather than as one possible outcome of a unified routing
decision.

\noindent\textbf{Vietnamese NLP and Legal QA.}
VnCoreNLP provides Vietnamese word segmentation and
NER~\cite{vu2018vncorenlp}, and PhoBERT provides pretrained Vietnamese
representations~\cite{nguyen2020phobert}. Legal NLP more broadly requires
reasoning over regulatory
relations~\cite{chalkidis2020legalbert,zhong2020nlp_law,guha2023legalbench,louis2024interpretable}.
This body of work supplies the linguistic and legal-domain foundations our
system builds on, but does not address adaptive retrieval routing.

In summary, despite substantial progress along each of these lines, no prior
system decides among dense retrieval, graph traversal, hybrid reasoning, and
clarification \emph{before} retrieval in a single framework. Advanced RAG
methods improve a single mechanism; GraphRAG methods assume traversal is
always needed; cost-aware routers select models rather than evidence
mechanisms; and ambiguity-handling methods treat clarification in isolation.
Our work differs in placing an explicit, reasoning-aware routing decision at
the centre of the pipeline, and in treating clarification as a first-class
route backed by deterministic conversational context resolution.
